Read in Gauge Theory Gravity, curvatures are linear maps on physical bivectors, such as the curvature operator $\mathcal{R}(B)$, and the orthonormal coframe and curvature two-forms named below are their components, a labelled reading (Fact \claimref{synthesis.imported-gtg-dictionary}, clauses (i) and (iv)); the closed Euclidean sections of (iii) lie beyond one flat chart. \emph{(i).} For linear maps $X, Y$ on bivectors write $E[X, Y] := 2\langle X(a\wedge b)\,I^{-1}\,Y(I\partial_b\wedge\partial_a)\rangle$, summed over a frame and its reciprocal, and $E[X] = E[X, X]$. In components, with $X$ read as the bivector-valued two-form $X^{ab} = \tfrac12 X^{ab}{}_{cd}\,e^{c}\wedge e^{d}$ in an orthonormal coframe and $\epsilon_{abcd}$ the permutation symbol with $\epsilon_{0123} = 1$, $e^{a}\wedge e^{b}\wedge e^{c}\wedge e^{d} = \epsilon_{abcd}\,\mathrm{vol}$ and the product of two permutation symbols is the generalised Kronecker delta in either signature, so $\epsilon_{abcd}X^{ab}\wedge Y^{cd} = \tfrac14\,\delta^{efgh}_{abcd}X^{ab}{}_{ef}Y^{cd}{}_{gh}\,\mathrm{vol} = E[X, Y]\,\mathrm{vol}$; the factor $I^{-1}$, which is $-I$ in Lorentzian and $+I$ in Euclidean signature, makes this hold in both. The three standard contractions are $E[\mathcal{R}] = E_{4}$, $E[\mathcal{R}, \mathrm{id}] = E[\mathrm{id}, \mathcal{R}] = 2R$ and $E[\mathrm{id}] = 24$ (in components $\epsilon_{abcd}R^{ab}\wedge R^{cd} = E_{4}\,\mathrm{vol}$, $\epsilon_{abcd}e^{a}\wedge e^{b}\wedge R^{cd} = 2R\,\mathrm{vol}$, $\epsilon_{abcd}e^{a}\wedge e^{b}\wedge e^{c}\wedge e^{d} = 24\,\mathrm{vol}$). On the model of constant curvature $s\ell^{-2}$ (de Sitter for $s = 1$, anti-de Sitter for $s = -1$) one has $\mathcal{R}(B) = s\ell^{-2}B$, so the Cartan curvature that vanishes there is $\mathcal{F} = \mathcal{R} - s\ell^{-2}\,\mathrm{id}$ (in components $F^{ab} = R^{ab} - s\ell^{-2}e^{a}\wedge e^{b}$); this is the sign convention of Fact \claimref{dark-energy.imported-macdowell-mansouri-theorem-stelle-west} after its curvature convention is matched, the content being that $F = 0$ on the model, as Theorem \claimref{dark-energy.frozen-point-makes-stelle-west-vector}(ii) states for the reduced connection. Expanding by bilinearity,
\[
\ell^{2}E[\mathcal{F}] = \ell^{2}E_{4} - 4sR + 24\ell^{-2} = -4s\Bigl(R - \frac{6s}{\ell^{2}} - \frac{s\ell^{2}}{4}E_{4}\Bigr),
\]
and with $\Lambda = 3s/\ell^{2}$ the bracket is $R - 2\Lambda + \alpha E_{4}$ with $\alpha = -s\ell^{2}/4 = -3/4\Lambda$. Normalising the Einstein-Cartan term to $R/16\pi G$ gives the Lagrangian $(R - 2\Lambda + \alpha E_{4})/16\pi G$. The coefficient is also fixed without the expansion: the square vanishes on the model, where $R = 4\Lambda$ and $E_{4} = 8\Lambda^{2}/3$, so $2\Lambda + \alpha\,8\Lambda^{2}/3 = 0$. The check verifies the expansion exactly on constant-curvature, product and generic pair-curvature tensors in both signatures and on both branches, reading $\epsilon_{abcd}$ from products of Clifford generators.

\emph{(ii).} Under (W), for a Lagrangian $(R - 2\Lambda + \alpha E_{4})/16\pi G$ the Wald entropy of a bifurcation surface is the Jacobson-Myers functional $\frac{1}{4G}\int_{\Sigma}(1 + 2\alpha\,{}^{(2)}R)\,dA$, ${}^{(2)}R$ the scalar curvature of the induced metric (clause (c) of Fact \claimref{synthesis.imported-wald-noether-charge-entropy}; (W) is the reading that this Wald entropy is the horizon's entropy). In two dimensions ${}^{(2)}R = 2K$, and for a closed surface $\int_{\Sigma}K\,dA = 2\pi\chi(\Sigma)$ (Gauss-Bonnet, Fact \claimref{synthesis.imported-chern-gauss-bonnet}(a)), so the Euler part is $\frac{2\alpha}{4G}\cdot 4\pi\chi = 2\pi\alpha\chi/G = -3\pi\chi/2\Lambda G$, a topological invariant, constant under every deformation of a closed cut. For an open cut, write the induced metric locally as $\varrho(dx^{2} + dy^{2})$ (a component reading in isothermal labels; $\varrho$ is a conformal factor, not the rotation-gauge field); then $\sqrt{h}\,{}^{(2)}R = -\Delta\ln\varrho$ with $\Delta$ the flat Laplacian, an identity the check verifies from the Christoffel symbols for $\varrho(1 + \varepsilon\psi)$ as rational functions of $(x, y, \varepsilon)$. A deformation $\varrho \mapsto \varrho(1 + \varepsilon\psi)$ changes the integrand by $-\Delta\ln(1 + \varepsilon\psi)$, a divergence whose flux vanishes when $\psi$ and its gradient vanish at the edge; at first order the change is $-\int\Delta\psi = 0$ (the check integrates $\Delta\psi$ exactly for $\psi = (1 - x^{2})^{2}(1 - y^{2})^{2}$ on the square). Since every two-dimensional metric is locally conformally flat (isothermal coordinates) and a diffeomorphism changes a scalar density by a divergence, the first variation of the Euler part is zero for every deformation supported away from the edge. When the deformation does not vanish with its gradient at the edge, the divergence leaves an edge flux: for $\psi = (1 - x^{2})(1 - y^{2})$, which vanishes on the edge but whose gradient does not, $\int\Delta\psi = -\tfrac{32}{3} \neq 0$ (in the check), so on open patches the statement is restricted as in the formal text. Any relation that compares $S$ on two closed cuts of one topology therefore involves only $\Delta A/4G$; and the coefficient $1/4G$ of the area is the Einstein-Cartan term's, a term of the series' square by (i). In the bulk the Euler term is exact (Theorem \claimref{dark-energy.wall-counterterm}(i) gives its primitive), consistent with its absence from the field equations.

\emph{(iii).} Take units of $\pi/\Lambda G$. De Sitter space has horizon radius $\sqrt{3/\Lambda}$, area $12\pi/\Lambda$, so $A/4G = 3$ (Fact \claimref{synthesis.imported-gibbons-hawking-entropy}; this is $\pi N^{2}$ of Proposition \claimref{synthesis.de-sitter-entropy-pi-n-squared}), and the shift for $\chi = 2$ is $-3$: total $0$. The Nariai solution is $\mathrm{dS}_{2}\times S^{2}$ with both radii $1/\sqrt{\Lambda}$, whose two horizons each have area $4\pi/\Lambda$: Einstein-Hilbert gives $1 + 1 = 2$ and the series' action $2 - 6 = -4$. For a Schwarzschild-de Sitter hole of mass $M$ with two horizons ($27M^{2}\Lambda < 1$), $h(r) = r f(r) = r - 2M - \Lambda r^{3}/3$ has $h(0) = -2M < 0$ and $h'(r) = 1 - \Lambda r^{2} > 0$ on $(0, 1/\sqrt{\Lambda})$, so its smaller positive root, the black-hole horizon, lies below $1/\sqrt{\Lambda}$; its area is below $4\pi/\Lambda$ and its entropy under the series' action below $1 - 3 = -2$. Euclidean check: with $I = -\frac{1}{16\pi G}\int(R - 2\Lambda + \alpha E_{4})\sqrt{g}$, the sphere $S^{4}$ of radius $\sqrt{3/\Lambda}$ has $R = 4\Lambda$, volume $8\pi^{2}\cdot 9/(3\Lambda^{2})$ and $\int E_{4} = 64\pi^{2} = 32\pi^{2}\chi(S^{4})$ (Fact \claimref{synthesis.imported-chern-gauss-bonnet}(b)), giving $I_{\mathrm{EH}} = -3$ and $I_{E_{4}} = +3$; the product $S^{2}\times S^{2}$ of radii $1/\sqrt{\Lambda}$ has $R = 4\Lambda$, volume $16\pi^{2}/\Lambda^{2}$ and $\int E_{4} = 128\pi^{2} = 32\pi^{2}\chi(S^{2}\times S^{2})$, giving $I_{\mathrm{EH}} = -2$ and $I_{E_{4}} = +6$. In each case $I = -\sum S$ over the horizons, the Euclidean $\chi$ being the sum of the horizons' $\chi$ (two fixed two-spheres in $S^{2}\times S^{2}$, one in $S^{4}$). The nucleation exponent $S_{\mathrm{Nariai}} - S_{\mathrm{dS}} = -(I_{\mathrm{Nariai}} - I_{\mathrm{dS}})$ is $2 - 3 = -1$ for Einstein-Hilbert and $-4 - 0 = -4$ for the series' action. For $\Lambda < 0$, $\alpha = -3/4\Lambda > 0$ and the shift is positive. All arithmetic here is exact and is reproduced in the check (with the black-hole root located by rational bisection at three masses).
